\section{Experimental Workflow and Evidence Production}
\label{sec:experiment-workflow}

GAWorld treats an experiment as a configured execution that produces an
inspectable artifact tree, rather than as a prompt followed by a single model
response.  The ordinary entry point, \texttt{generative\_city\_sim.py run},
loads the selected agents and city map, restores persistent state when enabled,
assembles the currently integrated plugins, and advances a shared simulation
clock.  Runtime configuration controls the horizon, agent cohort, random seed,
provider routing, environment, policy events, and output locations.  A run can
therefore be reconstructed at the level of its declared configuration and
stored traces, although exact replay also depends on external model behavior,
network services, and the completeness of the archived inputs.

\subsection{Matched Event and Baseline Runs}

For intervention studies, the \texttt{compare-event} command constructs two
scenario directories under a timestamped comparison root.  The
\texttt{without\_event} arm inherits the configured policy-event list, whereas
the \texttt{with\_event} arm appends a named event at an explicit simulation
day and clock time.  Both arms receive the same requested agent subset,
simulation horizon, random seed, and---when supplied---LLM provider override.
Before execution, each arm is reset within its own isolated memory, log,
vector-database, state, network, environment, intervention, and visualization
directories.  The two runs are then launched concurrently.  This organization
reduces accidental cross-arm state leakage and makes missing or failed arms
visible through separate reset and run logs.

After both subprocesses finish, the comparison routine reads each arm's
\path{state/agent_state_history.csv}.  For every available metric, it
computes the across-agent final snapshot, the mean over recorded observations,
and event-minus-baseline differences.  Machine-readable results are stored in
\path{comparison_metrics.csv}; \path{comparison_summary.md} provides a
human-readable view.  The separate \path{run_meta.json} records choices such
as seed, horizon, provider, and whether the lower-fidelity fast mode was
requested.  Policy-oriented fields can include stance, toxicity,
misinformation risk, cross-viewpoint exposure, and intervention reward when
their upstream mechanisms emit them.  Importantly, archived policy runs contain
fields that remain identically zero, including misinformation-risk and
cross-viewpoint-exposure series in some comparisons.  Presence in the schema
must therefore not be reported as evidence that the corresponding mechanism was
active.

\begin{figure}[t]
  \centering
  \includegraphics[width=\linewidth]{experiment_flow.pdf}
  \caption{GAWorld's experiment path separates a research question and
  configured scenario from matched execution arms, stored traces, metric
  extraction, and a descriptive audit.  The final boundary is substantive:
  an inspectable comparison is not by itself a real-world causal estimate.}
  \label{fig:experiment-flow}
\end{figure}

\subsection{What Matching Does---and Does Not---Identify}

Sharing a pseudorandom seed and configuration is a software control, not a
causal-identification strategy.  LLM endpoints can be nondeterministic despite
fixed sampling parameters; concurrent calls may experience different failures,
fallbacks, latency, or service-side model revisions; and an injected event can
change subsequent prompts and thereby amplify trajectory divergence.  Moreover,
the implemented report contrasts raw trajectories without random assignment,
replication, uncertainty estimation, or adjustment to an external population.
Consequently, GAWorld's paired output supports debugging, mechanism tracing,
and hypothesis generation.  It does not establish that a measured difference
was caused by the analogous intervention in a real city, or even that a
single-run simulation contrast is stable across seeds.

A defensible study should therefore pre-specify outcomes and event timing,
archive the exact configuration and provider information, repeat both arms over
multiple seeds, include negative controls where appropriate, and report
missingness and failures alongside effect summaries.  The event should also be
traced to the subsystem expected to implement it: textual exposure to a policy
is not equivalent to a verified economic, mobility, or network-state update.
This distinction is especially important in an LLM-based society, where a
plausible narrative response can mask a disconnected structural mechanism.

\subsection{Heterogeneous but Inspectable Outputs}

The current runtime writes several complementary views of a simulation.  Agent
logs preserve local action and cognition traces; state histories provide a long
table of agent, step, metric, and value; environment output contains the shared
timeline and contextual events; memory directories contain durable episodic,
relationship, growth, and vector-retrieval state; and the economy subsystem
maintains financial snapshots and audit artifacts.  Visualization traces encode
frames for replay, while intervention plugins retain policy-specific metrics.
Comparison roots place two complete scenario trees beside their derived
summary files.  Together these artifacts allow an analyst to move from an
aggregate difference back to the relevant agents, times, and subsystem records.

This plurality is currently both a strength and a limitation.  The emerging
kernel provides an append-only \texttt{Recorder} that time-stamps plugin events
into named JSONL streams, but it has not replaced or normalized the simulator's
existing logs, CSV files, environment timeline, memories, and subsystem-owned
outputs.  Cross-subsystem analysis still requires explicit joins and attention
to distinct schemas and clocks.  We therefore describe GAWorld as inspectable,
not as having a fully unified provenance store.

\subsection{GAWorld-Bench as an Audit Layer}

GAWorld-Bench consumes existing output trees or, when explicitly requested,
orchestrates comparison runs.  Its implemented tracks examine macro-level
anchors and selected intervention diagnostics, including known-sign checks,
placebo handling, and an optional determinism comparison.  The scorecard marks
unassessed components rather than treating missing evidence as success, and
labels its composite as a diagnostic summary rather than a scientific result.
Other proposed validation dimensions remain unimplemented or incomplete.
Synthetic fixtures bundled with the harness exercise parsers and scoring paths;
they are software tests and are never treated in this paper as observations from
an artificial society.  The benchmark thus complements the artifact workflow
in \cref{fig:experiment-flow}: it helps expose coverage gaps and suspicious
outputs, but cannot repair an underpowered or poorly identified experiment.
